%% ============================================================================
%%  Stochastic ODE-Guard: A Neural Stochastic Differential Equation Framework
%%  with Anti-Concentration Bounds for Adversarially Robust Network Intrusion
%%  Detection
%%
%%  Version v4 — IEEE Access resubmission (revised after peer review).
%%  Target venue : IEEE Access
%%  Authors      : R. N. Anaedevha, A. G. Trofimov, Y. V. Borodachev
%%  Affiliation  : National Research Nuclear University MEPhI, Moscow, Russia
%%  Compile with : pdflatex on Overleaf using the IEEE Access template
%%                 (ieeeaccess.cls is bundled with that Overleaf template)
%%
%%  Every change from v3 → v4 is intended to address a specific reviewer
%%  concern from the IEEE Access decision letter (2026-08). The full
%%  point-by-point response is in
%%  docs/response_to_reviewers/RESPONSE.md in the reproducibility repo.
%% ============================================================================

\documentclass{ieeeaccess}

%%%% PACKAGES
\usepackage[nocompress]{cite}
\usepackage{amsmath,amssymb,amsfonts,amsthm,mathtools}
\usepackage{bm}
\usepackage{graphicx}
\graphicspath{{./}{figures/}{stochastic_odeguard_figures/}}
\usepackage{textcomp}
\usepackage{booktabs}
\usepackage{multirow}
\usepackage{array}
\usepackage{algorithm}
\usepackage{algorithmic}
\usepackage{subcaption}
\usepackage{enumitem}
\usepackage[hyphens]{url}
\usepackage[hidelinks]{hyperref}

%%%% Theorem environments
\newtheorem{theorem}{Theorem}
\newtheorem{lemma}[theorem]{Lemma}
\newtheorem{proposition}[theorem]{Proposition}
\newtheorem{corollary}[theorem]{Corollary}
\newtheorem{definition}[theorem]{Definition}
\newtheorem{assumption}[theorem]{Assumption}
\newtheorem{remark}[theorem]{Remark}

%%%% Math shortcuts
\newcommand{\Real}{\mathbb{R}}
\newcommand{\Exp}{\mathbb{E}}
\newcommand{\Prob}{\mathbb{P}}
\newcommand{\Var}{\mathrm{Var}}
\newcommand{\KL}{\mathrm{KL}}
\newcommand{\Lip}{\mathrm{Lip}}
\newcommand{\Brn}{\mathrm{B}}
\newcommand{\Fcal}{\mathcal{F}}
\newcommand{\Gcal}{\mathcal{G}}
\newcommand{\Xcal}{\mathcal{X}}
\newcommand{\Ycal}{\mathcal{Y}}
\newcommand{\Hcal}{\mathcal{H}}
\newcommand{\Ncal}{\mathcal{N}}
\newcommand{\Acal}{\mathcal{A}}
\newcommand{\Dcal}{\mathcal{D}}
\newcommand{\Lcal}{\mathcal{L}}
\newcommand{\Zcal}{\mathcal{Z}}
\newcommand{\pinv}{{+}}   % Moore--Penrose pseudo-inverse superscript

\hyphenation{op-tical net-works semi-conduc-tor}
\hypersetup{pdfauthor={Roger Nick Anaedevha}, pdftitle={Stochastic ODE-Guard: A Neural Stochastic Differential Equation Framework with Anti-Concentration Bounds for Adversarially Robust Network Intrusion Detection}}

\begin{document}

\history{Date of publication xxxx 00, 0000, date of current version xxxx 00, 0000.}
\doi{10.1109/ACCESS.2026.DOI}

\title{Stochastic ODE-Guard: A Neural Stochastic Differential Equation Framework with Anti-Concentration Bounds for Adversarially Robust Network Intrusion Detection}

\author{Roger Nick Anaedevha\authorrefmark{1},
Alexander Gennadevich Trofimov\authorrefmark{1},
and Yuri Vladimirovich Borodachev\authorrefmark{2}}

\address[1]{National Research Nuclear University MEPhI (Moscow Engineering Physics Institute), Moscow 115409, Russia (e-mail: ar006@campus.mephi.ru)}
\address[2]{Artificial Intelligence Research Center, National Research Nuclear University MEPhI, Moscow 115409, Russia}

\tfootnote{This work was supported by a grant for research centers in Artificial Intelligence from the Ministry of Economic Development of the Russian Federation under subsidy agreement (identifier 000000C313925P3Q0002). All authors of this publication are co-first authors.}

\markboth{Anaedevha \headeretal: Stochastic ODE-Guard: A Neural Stochastic Differential Equation Framework for Adversarially Robust Network Intrusion Detection}{Anaedevha \headeretal: Stochastic ODE-Guard: A Neural Stochastic Differential Equation Framework for Adversarially Robust Network Intrusion Detection}

\corresp{Corresponding author: Roger Nick Anaedevha (e-mail: ar006@campus.mephi.ru).}

\begin{abstract}
Adversarially robust network intrusion detection remains an open problem
because per-flow neural classifiers degrade catastrophically under small
feature-space perturbations that lie within the empirical noise floor of
production flow extractors. Existing certified defences such as
Lipschitz-bounded networks and randomised smoothing trade clean
accuracy against certificate radius and rarely transfer to the
long-tailed, high-dimensional flow-feature spaces that production
intrusion detection and prevention systems must handle. We introduce
\textbf{Stochastic ODE-Guard}, a neural stochastic differential equation
(NSDE) framework that treats the embedded flow representation as the
terminal state of an It\^o diffusion whose drift and diffusion
coefficients are jointly learned. The framework admits three
theoretical guarantees: an \emph{empirically certified} $L^{2}$
Lipschitz bound on the smoothed margin (Lemma~\ref{lem:closed-lg}), a
Bismut--Elworthy--Li identity with a Moore--Penrose pseudo-inverse
diffusion that yields stable gradient estimators for non-square
diffusion parameters (Theorem~\ref{thm:bel}), and---as our principal
contribution---a \emph{probabilistic robustness proposition}
(Proposition~\ref{prop:robust}) that follows from the Carbery--Wright
anti-concentration inequality applied to the Wiener-chaos expansion of
the smoothed score. The chaos truncation degree is now selected
adaptively per example via Hermite regression rather than fixed a
priori, and the certificate uses a Hoeffding \emph{lower} confidence
bound on the margin $L^{2}$ norm in the Carbery--Wright denominator,
which restores the direction of the inequality. We rename the resulting
quantity a \emph{probabilistic robustness radius} to distinguish it
from deterministic $\arg\max$ certificates; the corresponding
$\arg\max$ decision-flip bound is derived by a union argument
(Corollary~\ref{cor:argmax}). Empirical validation uses three
Kaggle-released benchmark suites (ICS3D, IIS3D, IDS-PQC) totalling
$1.84{\times}10^{7}$ flows, with new provenance, deduplication and
leakage-analysis appendices; the framework is evaluated under random,
temporal, host-disjoint and scenario-disjoint splits, and against
threat-model-matched neural baselines with commercial signature engines
and LLM-safety filters reported separately as \emph{reference
baselines}. Stochastic ODE-Guard achieves $96.4\%$ clean F1 and
$93.1\%$ F1 under PGD-40 at $\epsilon{=}0.03$ (random splits),
$91.8\%$ under the temporal-holdout protocol, and $90.6\%$ under the
host-disjoint protocol; the strongest deployed baseline (SDE-TGNN)
reaches $87.2\%$, $84.5\%$, and $82.8\%$ respectively. Under the
problem-space projection that enforces packet-grammar feasibility,
SODE-Guard's lead grows to $+7.2$ F1 points at $\epsilon{=}0.03$.
Median inference latency is $1.7$~ms on a single A100 GPU.
\end{abstract}

\begin{keywords}
Adversarial robustness, anti-concentration inequalities, network intrusion
detection, neural stochastic differential equations, probabilistic
robustness radius, uncertainty quantification.
\end{keywords}

\titlepgskip=-21pt

\maketitle

\section{Introduction}\label{sec:intro}
%% =====================================================================
\IEEEPARstart{N}{etwork} intrusion detection and prevention systems
(IDPS) have become the principal line of defence against the escalating
volume of cyber-physical threats targeting cloud, edge, and
Internet-of-Things deployments. Modern neural IDPS deliver clean-traffic
F1 scores exceeding $0.97$ on canonical benchmarks such as UNSW-NB15
and CIC-IDS2017~\cite{wu2022rtids,halbouni2022cnnlstm}, yet the same
models routinely lose more than fifteen F1 points under
projected-gradient adversarial perturbations whose $\ell_\infty$ budget
is smaller than the empirical noise floor of the underlying flow
extractor~\cite{ennaji2024advsurvey,vitorino2024advbench}. The
vulnerability is not merely academic: structural adversarial attacks on
graph-neural-network IDPS have been demonstrated to evade
production-grade detectors at line rate~\cite{venturi2024structural}.

The dominant certified-robustness strategies---Lipschitz-bounded
networks~\cite{tsuzuku2018lipschitz,anil2019sorting,araujo2023sll}
and randomised smoothing~\cite{cohen2019smoothing,salman2019smoothadv,yang2020smoothing}---deliver
formal guarantees but suffer two structural limitations when
transplanted into IDPS. First, Lipschitz certificates degrade linearly
with input dimension, which is punishing for the $83$-dimensional flow
vectors produced by modern extractors~\cite{anil2019sorting}. Second,
randomised smoothing inflates inference latency by the Monte-Carlo
sample budget required for tight Clopper--Pearson certificates, which
is incompatible with the sub-millisecond per-flow envelope of production
sensors~\cite{cohen2019smoothing}. Both families also assume the
classifier is a deterministic map, ignoring the natural stochasticity
of continuous-time network dynamics.

A complementary line of work has shown that continuous-time neural
dynamics---neural ordinary differential equations
(NODEs)~\cite{kang2021sodef,guglielmi2025minimal}---inherit beneficial
robustness from the Gr\"onwall--Bellman inequality on their flow map.
Extending the framework from NODEs to neural stochastic differential
equations (NSDEs)~\cite{tzen2019nsde,li2020scalable,kidger2021nsdegan}
introduces an explicit diffusion term whose role under adversarial
perturbation has, until now, been left implicit. Recently the
\emph{stable Neural SDE} family~\cite{oh2024stablesde} established
existence and uniqueness conditions on the irregular-time-series
setting, and applications to graph uncertainty~\cite{bergna2024lgnsde}
confirm that the diffusion channel carries calibrated epistemic
information. What has been missing is a tractable robustness theory
that links the SDE diffusion coefficient to a finite-sample
adversarial certificate.

This paper closes that gap. We introduce \textbf{Stochastic ODE-Guard},
a neural stochastic differential equation framework whose drift and
diffusion are jointly learned on the embedded flow representation,
whose gradients are estimated through a Bismut--Elworthy--Li identity
with a Moore--Penrose pseudo-inverse~\cite{bismut1981martingales,elworthy1994formulae},
and whose adversarial robustness derives from an anti-concentration
argument on the Wiener-chaos expansion of the smoothed score. The
framework was prototyped inside the deployed \texttt{RobustIDPS.ai}
open-source platform whose registry of fourteen detectors provides
both infrastructure and threat-model-matched competitive baselines,
and was validated on three Kaggle-released benchmark suites totalling
$1.84\times 10^7$ flows.

\medskip\noindent\textbf{What is genuinely new in the composition.}
Neural SDEs, stochastic smoothing, PAC-Bayes bounds and NIDS
benchmarking are all established individually; our contribution is
their principled composition together with three technical results
that, to the best of our knowledge, are new. Specifically:

\emph{(C1)~The anti-concentration probabilistic robustness proposition.}
Section~\ref{sec:theory} proves that the smoothed-score gap between a
clean input and an adversarial perturbation is anti-concentrated in
the sense of Carbery and Wright~\cite{carbery2001anticonc}. The bound
is polynomial in the perturbation budget and identifies the effective
Wiener-chaos degree $d^\star$ as the controlling parameter. We
rebuilt this proposition in the present revision so that
Carbery--Wright is applied with a Hoeffding \emph{lower} confidence
bound on the polynomial $L^{2}$ norm rather than an upper bound
(Section~\ref{sec:theory}; Appendix~B). To our knowledge this is the
first robustness certificate that follows from the algebraic structure
of the diffusion rather than from Lipschitz contraction or
Gaussian-noise smoothing, and the corresponding
$\arg\max$-decision-flip bound is derived by a union argument in
Corollary~\ref{cor:argmax}.

\emph{(C2)~A Bismut--Elworthy--Li estimator for a non-square diffusion.}
Section~\ref{sec:framework} derives the gradient of the expected loss
with respect to the diffusion coefficient using the Bismut--Elworthy--Li
formula~\cite{bismut1981martingales,elworthy1994formulae,hsu2002}
adapted to the $d\times m$ diffusion setting with $m<d$. The novelty is
the constructive use of the Moore--Penrose pseudo-inverse
$g_\theta^{\pinv} = (g_\theta^\top g_\theta + \lambda_0 I_m)^{-1}
g_\theta^\top$, whose bounded condition number is guaranteed by the
ellipticity floor $\lambda_0 = 10^{-3}$. Appendix~D reports its
empirical bias and variance against a finite-difference reference.

\emph{(C3)~An integrated, split-robust ablation across three Kaggle
benchmarks and the deployed platform.} Section~\ref{sec:exp} reports
the first unified evaluation of Stochastic ODE-Guard against six
threat-model-matched neural baselines and three commercial /
LLM-policy reference baselines on ICS3D~\cite{ds_ics3d},
IIS3D~\cite{ds_iis3d}, and the encrypted post-quantum-cryptography
traffic corpus~\cite{ds_pqcids}. Every experiment is repeated under
four split protocols (random, temporal, host-disjoint,
scenario-disjoint) and under two threat models (feature-space and
problem-space with packet-grammar feasibility projection).

\emph{(C4)~A production-grade reproducibility artefact.} The entire
codebase, including the new $L^{2}$-Lipschitz certifier, the adaptive
chaos-degree estimator, the pseudo-inverse-based BEL estimator, the
PAC-Bayes-kl certificate, the four split protocols, the feasibility
projector, and the sensitivity-swept randomised-smoothing baseline, is
released under MIT license and is described module-by-module in the
reproducibility index at
\url{https://github.com/rogerpanel/CV/tree/main/SODE-Guard}.

The remainder is organised as follows. Section~\ref{sec:related}
surveys related work in four threads. Section~\ref{sec:prelim} states
the mathematical preliminaries. Section~\ref{sec:framework} presents
the Stochastic ODE-Guard architecture and training procedure.
Section~\ref{sec:theory} proves the principal anti-concentration
robustness result together with its arg-max corollary.
Section~\ref{sec:exp} reports the experimental setup, the integrated
ablation, and the industry comparison. Section~\ref{sec:disc}
discusses limitations and threats to validity, and
Section~\ref{sec:conc} concludes.

\medskip\noindent\textbf{Intuition.} A network flow is a sequence of
packet-level statistics that a deployed sensor extracts in real time.
A classical neural classifier treats this flow as a fixed vector and
predicts a label deterministically. Stochastic ODE-Guard does not
perturb the input; instead, it \emph{embeds} the flow into a
$128$-dimensional feature representation and evolves that embedded
representation forward in continuous time under a learned It\^o
diffusion. The final classification averages the softmax head over a
small number of Brownian samples of the embedded trajectory. The
diffusion is what gives the framework its adversarial robustness,
because a small perturbation of the input has limited ability to
change the \emph{averaged} terminal state when the average is over a
non-degenerate Brownian ensemble. The theoretical contribution
quantifies exactly how limited that ability is.

%% =====================================================================
\section{Related Work}\label{sec:related}
%% =====================================================================
Four bodies of literature converge on this work.

\emph{Certified adversarial robustness.} The Lipschitz-bounded family
established by Tsuzuku \textit{et~al.}~\cite{tsuzuku2018lipschitz}
and refined by Anil \textit{et~al.}'s
GroupSort~\cite{anil2019sorting}, Wang and Manchester's
Lipschitz-bounded deep networks~\cite{wang2023lbdn}, and Araujo
\textit{et~al.}'s SLL parameterisation~\cite{araujo2023sll}, provides
certificates whose radius scales as the inverse network-Lipschitz
constant. Randomised smoothing~\cite{cohen2019smoothing} and its
anisotropic extensions~\cite{yang2020smoothing,eiras2022ancer} convert
non-robust base classifiers into smoothed classifiers with
high-probability $\ell_2$ certificates. PixelDP~\cite{lecuyer2019pixeldp}
established the differential-privacy connection that grounds the
Gaussian-smoothing analysis. Section~\ref{sec:theory} contributes a
certificate that is neither Lipschitz nor smoothing-based: it exploits
the anti-concentration of the polynomial-chaos representation induced
by the SDE diffusion. The closest neighbour theoretically is the
Carbery--Wright bound~\cite{carbery2001anticonc} as developed by Meka,
Nguyen, and Vu~\cite{meka2016anticonc}, but neither was previously
applied to deep-learning robustness. In particular, our derivation
uses the exact $L^{2}$ norm of the chaos-truncated polynomial as the
Carbery--Wright denominator (Appendix~B), replacing an upper bound
that would reverse the inequality's direction.

\emph{Neural ODEs and SDEs for robustness.} Kang \textit{et~al.}'s
SODEF~\cite{kang2021sodef} showed that Lyapunov-stable equilibrium
points in a neural ODE confer adversarial robustness. Guglielmi
\textit{et~al.}~\cite{guglielmi2025minimal} quantified the minimal
weight perturbation under which a neural ODE loses robustness.
Stable Neural SDEs~\cite{oh2024stablesde} established well-posed
extensions to irregular time series, while Bergna \textit{et~al.}'s
latent graph neural SDE~\cite{bergna2024lgnsde} carried these ideas
into the graph setting. The stochastic-adjoint sensitivity method of
Li \textit{et~al.}~\cite{li2020scalable} provided the practical
training algorithm. Our work treats the diffusion as a robustness
mechanism rather than a generative-modelling tool (the latter is the
focus of score-based SDEs~\cite{song2021score}) and derives a
non-trivial certificate from its algebraic structure, extended
constructively to a non-square diffusion coefficient via the
Moore--Penrose pseudo-inverse.

\emph{Network intrusion detection.} Lo \textit{et~al.}'s
E-GraphSAGE~\cite{lo2022egraphsage} brought edge-feature graph neural
networks to NIDS. Wu \textit{et~al.}'s RTIDS~\cite{wu2022rtids} adapted
the transformer to encrypted-flow classification. Halbouni
\textit{et~al.}'s hybrid CNN-LSTM~\cite{halbouni2022cnnlstm}
established the spatio-temporal template subsequently refined by
Mosaiyebzadeh \textit{et~al.}'s
IDS-GraphMamba~\cite{mosaiyebzadeh2025graphmamba}. Our previous work
contributed three building blocks to this line: a temporal-adaptive
neural ODE with deep spatio-temporal point processes for real-time
NIDS~\cite{anaedevha2026ode}, a stochastic multimodal transformer
with uncertainty quantification~\cite{anaedevha2025stochtx}, and a
robust-adversarial defence formulation~\cite{anaedevha2024ram}. In
addition we examined device-identification baselines on wireless
intrusion traffic~\cite{anaedevha2024elcon}.

\emph{Stochastic neural networks and uncertainty.} The Bayesian
neural network literature beginning with MacKay~\cite{mackay1992bnn}
and extended in the adversarial setting by Liu \textit{et~al.}'s
Adv-BNN~\cite{liu2019advbnn} demonstrated that posterior weight
uncertainty confers some adversarial robustness, an effect quantified
by He \textit{et~al.}'s parametric noise injection~\cite{he2019pni}.
The PAC-Bayes school~\cite{mcallester1999pacbayes,dziugaite2017pacbayes,
perezortiz2021pacbayes} provides generalisation certificates for
stochastic networks that we use in Section~\ref{sec:theory} as a
complement to the anti-concentration argument and connect explicitly
to the empirical calibration error in Section~\ref{subsec:pacbayes-emp}.
Our prior work on MambaShield extended selective state-space modelling
to the poisoning-resilient setting~\cite{anaedevha2026mambashield},
and our uncertainty-calibrated hierarchical Gaussian
processes~\cite{anaedevha2026uchgp} furnished the multi-scale temporal
calibration we adopt here in Section~\ref{subsec:training}.

%% =====================================================================
\section{Mathematical Preliminaries}\label{sec:prelim}
%% =====================================================================
Throughout, $\Real^d$ denotes the $d$-dimensional Euclidean space,
$W=(W_t)_{t\geq 0}$ denotes an $m$-dimensional standard Brownian motion
on a filtered probability space
$(\Omega,\Fcal,\{\Fcal_t\}_{t\geq 0},\Prob)$ satisfying the usual
conditions, $\|\cdot\|$ is the Euclidean norm, $I_d$ the $d\times d$
identity matrix, and $A\succeq B$ for symmetric matrices means $A-B$
is positive semi-definite.

\subsection{Neural Stochastic Differential Equations}\label{subsec:nsde}
A neural stochastic differential equation on $[0,T]$ is the It\^o SDE
\begin{equation}\label{eq:nsde}
  dX_t = f_\theta(X_t,t)\,dt + g_\theta(X_t,t)\,dW_t,\qquad X_0=x_0,
\end{equation}
with $f_\theta:\Real^d\!\times\![0,T]\!\to\!\Real^d$ and
$g_\theta:\Real^d\!\times\![0,T]\!\to\!\Real^{d\times m}$ neural
networks parameterised by $\theta$.

\begin{theorem}[Existence and uniqueness; {\protect\cite[Thm.~5.2.1]{oksendal2003}}]\label{thm:eu}
If the drift $f_\theta$ and diffusion $g_\theta$ are Borel-measurable
and satisfy, uniformly in $t\in[0,T]$, the linear-growth condition
$\|f_\theta(x,t)\|+\|g_\theta(x,t)\|\leq C(1+\|x\|)$
and the Lipschitz condition
$\|f_\theta(x,t)-f_\theta(y,t)\|+\|g_\theta(x,t)-g_\theta(y,t)\|\leq
C\|x-y\|$ for a constant $C>0$, and if $\Exp\|X_0\|^2<\infty$, then
\eqref{eq:nsde} admits a unique $\Fcal_t$-adapted strong solution.
\end{theorem}

In practice the linear-growth and Lipschitz conditions are enforced by
spectral normalisation of the neural-network weights of $f_\theta$ and
$g_\theta$ as in Anil \textit{et~al.}~\cite{anil2019sorting}; the
implementation-level audit is deferred to Assumption~\ref{ass:impl}.

\subsection{Bismut--Elworthy--Li Identity for Non-Square Diffusion}
\label{subsec:bel}

We adapt the classical BEL formula~\cite{bismut1981martingales,
elworthy1994formulae,hsu2002} to the non-square diffusion case
$g:\Real^d\!\to\!\Real^{d\times m}$ with $m\leq d$. Let $g^{\pinv}$
denote the Moore--Penrose pseudo-inverse
\begin{equation}\label{eq:pinv}
   g^{\pinv}(x) \;=\; \bigl(g(x)^\top g(x) + \lambda_0 I_m\bigr)^{\!-1}
                       g(x)^\top,
\end{equation}
regularised by the ellipticity floor $\lambda_0>0$ so that the
inversion in~\eqref{eq:pinv} is well-conditioned.

\begin{theorem}[BEL identity with pseudo-inverse]\label{thm:bel}
Assume $f\in C^2_b(\Real^d,\Real^d)$ and
$g\in C^2_b(\Real^d,\Real^{d\times m})$ with the regularised
Gram matrix $g^\top g + \lambda_0 I_m$ uniformly positive-definite
in $x$. Let $X^x$ solve $dX_t=f(X_t)\,dt+g(X_t)\,dW_t$ with
$X_0=x$ and let $J^x_t=\partial X^x_t/\partial x$ denote the
first-variation process. Then for every bounded measurable
$\phi:\Real^d\!\to\!\Real$ and every $t>0$,
\begin{equation}\label{eq:bel}
  \nabla_x\,\Exp[\phi(X^x_t)]
  = \Exp\!\left[\phi(X^x_t)\,\frac{1}{t}\!\int_0^t\!
      \!\bigl(g^{\pinv}(X^x_s)\,J^x_s\bigr)^{\!\top}\!dW_s\right].
\end{equation}
\end{theorem}

In~\eqref{eq:bel} the pseudo-inverse $g^{\pinv}$ replaces the (undefined
in the non-square setting) matrix inverse used in the classical
statement, and the ellipticity floor $\lambda_0$ bounds the condition
number of $g^\top g + \lambda_0 I_m$ (empirical median $8.9$, worst-case
$32.4$ in our runs; see Appendix~D). The right-hand side of
\eqref{eq:bel} contains no derivative of $\phi$ and therefore extends
to discontinuous safety thresholds; this is essential because the IDPS
decision boundary is a step function of the terminal-state classifier
score.

\subsection{Carbery--Wright Anti-Concentration}\label{subsec:cw}
\begin{theorem}[Carbery--Wright, 2001~\cite{carbery2001anticonc}]\label{thm:cw}
There exists an absolute constant $C>0$ such that for every
$n,d\in\mathbb{N}$, every polynomial $p:\Real^n\!\to\!\Real$ of degree
at most $d$, and every $\varepsilon>0$,
\begin{equation}\label{eq:cw}
  \Prob_{Z\sim\Ncal(0,I_n)}\!\left[\,|p(Z)|\le
   \varepsilon\,\bigl(\Exp\,p(Z)^2\bigr)^{1/2}\right]
   \leq C\,d\,\varepsilon^{1/d}.
\end{equation}
\end{theorem}

The denominator is the \emph{exact} $L^{2}$ norm of $p$. The revised
Proposition~\ref{prop:robust} uses a Hoeffding lower confidence bound
on this quantity, not an upper bound, so that the direction of the
inequality is preserved (see Appendix~B).

\subsection{PAC-Bayes Generalisation for Stochastic Predictors}
\label{subsec:pacbayes}
\begin{theorem}[PAC-Bayes-kl; Maurer~\cite{maurer2004pacbayes}]\label{thm:pacbayes}
Let $\ell:\Hcal\times\Zcal\!\to\![0,1]$ be a loss bounded in $[0,1]$,
let $P$ be a data-independent prior on $\Hcal$, let $S\sim\Dcal^n$ be
an i.i.d.\ sample, and let $\delta\in(0,1)$. With probability at least
$1-\delta$ over the draw of $S$, simultaneously for every posterior
$Q$,
\begin{equation}\label{eq:pacbayes}
 \mathrm{kl}\!\left(\widehat{R}_S(Q)\,\big\|\,R(Q)\right)
   \leq \frac{\KL(Q\|P)+\log(2\sqrt{n}/\delta)}{n},
\end{equation}
where $\widehat R_S(Q)=\Exp_{h\sim Q}\widehat R_S(h)$ and
$R(Q)=\Exp_{h\sim Q}R(h)$.
\end{theorem}

We use the kl-inversion of~\eqref{eq:pacbayes} (Reeb \& Seldin,
2014~\cite{reeb2014klinv}) implemented in
\texttt{src/theory/pac\_bayes.py} of the reproducibility repository.
Section~\ref{subsec:pacbayes-emp} pairs the returned bound with the
empirical calibration error for each benchmark.

%% =====================================================================
\section{The Stochastic ODE-Guard Framework}\label{sec:framework}
%% =====================================================================
Section~\ref{sec:prelim} furnished the theoretical building blocks;
we now assemble them into the operational framework. The architectural
pipeline is summarised in Fig.~\ref{fig:architecture} and detailed in
three subsections: flow representation (\ref{subsec:flowrep}), the
SDE-based dynamics (\ref{subsec:sde}), and the training procedure with
the BEL gradient estimator (\ref{subsec:training}).

\begin{assumption}[Analytic $\leftrightarrow$ implementation]\label{ass:impl}
Throughout, $f_\theta$ and $g_\theta$ are three-layer MLPs with GELU
activations ($C^{\infty}$ smooth). Each linear layer $W\in\Real^{h\times h'}$
is spectral-normalised so $\|W\|_{2}\leq 1$; the resulting drift and
diffusion Lipschitz constants are bounded by
$K_f, K_g \leq \prod_\ell \|W_\ell\|_2 \cdot \|\phi\|_\Lip \leq 1$
where $\|\phi\|_\Lip$ is the GELU Lipschitz constant. The diffusion is
projected at every integration step through
$g_\theta(x,t)\to g_\theta(x,t)+\sqrt{\lambda_0}\,I_{d\wedge m}$ with
$\lambda_0=10^{-3}$, so $g_\theta^\top g_\theta \succeq \lambda_0 I_m$
uniformly and $g_\theta^{\pinv}$ from~\eqref{eq:pinv} is well-defined.
\end{assumption}

\begin{figure*}[!t]
\centering
\includegraphics[width=\textwidth]{figures/fig1_SODEGuard_arch.png}
\caption{Stochastic ODE-Guard architecture (unchanged from v3, retained
for orientation). The forward path (solid arrows) ingests the three
Kaggle benchmark suites, extracts $83$-dimensional flow features and
edge representations, propagates the embedded state through the It\^o
SDE of \eqref{eq:nsde} whose drift $f_\theta$ and diffusion $g_\theta$
are jointly learned, and emits a $34$-class allow/block verdict at the
terminal time $T$. The theorem-backed gradient path (dashed) exposes
the two analytical guarantees that distinguish the framework: the
Bismut--Elworthy--Li identity of Theorem~\ref{thm:bel} supplies a
gradient estimator that is unbiased under Assumption~\ref{ass:impl}
and uses the Moore--Penrose pseudo-inverse of~\eqref{eq:pinv}, while
the Carbery--Wright anti-concentration argument
(Proposition~\ref{prop:robust}) certifies the adversarial robustness
of the smoothed classification.}
\label{fig:architecture}
\end{figure*}

\subsection{Flow Representation}\label{subsec:flowrep}
Each network flow is represented by a feature vector
$x\in\Real^{83}$ extracted from the deployed RobustIDPS.ai pipeline,
augmented with edge-feature embeddings computed by an
E-GraphSAGE-style aggregator~\cite{lo2022egraphsage} over a sliding
micro-graph of size $B$ flows. The graph view encodes inter-flow
relationships that bare per-flow vectors miss. Letting
$h(x)\in\Real^{128}$ denote the concatenation of the per-flow features
and the edge-aggregated embedding, we set $X_0=h(x)$ in the SDE
\eqref{eq:nsde}. \emph{The SDE evolves the embedded state $X_t$; the
raw input feature vector $x$ is never perturbed by the diffusion. Any
adversarial perturbation in this paper is applied to the input side
$x$, and the certificate quantifies its effect after propagation
through the network and the SDE.}

\subsection{SDE Dynamics}\label{subsec:sde}
The dynamics are governed by the It\^o SDE
\begin{equation}\label{eq:sde-train}
\begin{aligned}
  dX_t &= f_\theta(X_t,t)\,dt + g_\theta(X_t,t)\,dW_t,\\
  X_0  &= h(x),\quad t\in[0,1],
\end{aligned}
\end{equation}
where the drift $f_\theta:\Real^{128}\!\times\![0,1]\!\to\!\Real^{128}$
and diffusion $g_\theta:\Real^{128}\!\times\![0,1]\!\to\!\Real^{128\times 16}$
are realised as 3-layer MLPs (hidden width $256$, GELU) with spectral
normalisation (Assumption~\ref{ass:impl}). The diffusion image
dimension $m=16$ keeps the Brownian-motion cost low while enabling the
polynomial-chaos expansion of Section~\ref{sec:theory}. We use a
fixed-step Euler--Maruyama integrator with $\Delta t = 0.05$ (20 steps)
at inference and the virtual-Brownian-tree variance-reduction scheme
of Li \textit{et~al.}~\cite{li2020scalable} during training. Under
Assumption~\ref{ass:impl} the Milstein correction at $\Delta t/2$
introduces a median $1.4\times 10^{-4}$ discrepancy on the terminal
state (Appendix~E), justifying the Euler--Maruyama choice at inference.

The terminal state $X_T$ feeds a linear classification head
$\psi_\eta:\Real^{128}\!\to\!\Real^{34}$ whose output is the
$34$-dimensional logit vector. The predictor is
\begin{equation}\label{eq:predictor}
  \hat y(x) \;=\; \mathrm{softmax}\!\bigl(\Exp[\psi_\eta(X_T)\mid X_0=h(x)]\bigr),
\end{equation}
where the expectation is realised by Monte-Carlo averaging over
$N_{\mathrm{mc}}=8$ Brownian paths at inference. The choice
$N_{\mathrm{mc}}=8$ is calibrated in Section~\ref{sec:exp} to keep
latency below $2$\,ms while controlling the variance of the smoothed
score.

\subsection{Training Procedure}\label{subsec:training}
The optimisation objective combines cross-entropy
$\Lcal_{\mathrm{CE}}$ with an anti-concentration regulariser
$\Lcal_{\mathrm{AC}}$ that we now define precisely (Reviewer~1 concern).
Let $m_i(\theta) := \psi_\eta^{(1)}(X_T^{(i)}) - \psi_\eta^{(2)}(X_T^{(i)})$
denote the top-1 minus top-2 logit gap for example $i$ averaged over
$P$ Brownian paths, let
$\hat s_i := \bigl(\tfrac{1}{P}\sum_{p=1}^P m_i^{(p)}(\theta)^2\bigr)^{1/2}$
be its empirical $L^{2}$ norm, and let
$\sigma_\kappa(u) := (1+\exp(-\kappa u))^{-1}$ be a temperature-$\kappa$
sigmoid surrogate for the indicator. The regulariser is
\begin{equation}\label{eq:ac-loss}
  \Lcal_{\mathrm{AC}}(\theta) \;=\; \frac{1}{|\mathcal B|}\!\sum_{\beta\in\mathcal B}\!
   \log\!\Bigl(1 + C\,d^\star\,(\beta / \hat s)^{1/d^\star}\,
                 \tfrac{1}{B}\!\sum_{i=1}^B\!\sigma_\kappa(\beta-|m_i|)\Bigr),
\end{equation}
with $\mathcal B=\{0.01,0.025,0.05,0.10\}$, $\kappa=50$, and
$d^\star=4$ during training (the certificate at inference uses the
adaptive $d^\star_{p99}$ of Algorithm~\ref{alg:chaos}). The complete
composite loss is
\begin{equation}\label{eq:loss}
  \Lcal(\theta,\eta) \;=\; \Lcal_{\mathrm{CE}}(\theta,\eta) \;+\;
   \lambda\,\Lcal_{\mathrm{AC}}(\theta),
\end{equation}
with $\lambda=0.1$ chosen on the IIS3D validation split.

Gradients with respect to the drift $\theta_f$ are computed through the
stochastic-adjoint sensitivity method~\cite{li2020scalable}. Gradients
with respect to the diffusion $\theta_g$ use the BEL identity
\eqref{eq:bel} of Theorem~\ref{thm:bel}, instantiated as
\begin{equation}\label{eq:bel-grad}
  \nabla_{\theta_g}\,\Exp[\psi_\eta(X_T)]
   \,=\,\Exp\!\left[\psi_\eta(X_T)\,M_T(\theta_g)\right],
\end{equation}
where $M_T(\theta_g) = \tfrac{1}{T}\int_0^T
(g_\theta^{\pinv} J_s)^\top dW_s$ is the Malliavin weight with
$g_\theta^{\pinv}$ from \eqref{eq:pinv}. The Adam optimiser is run
with learning rate $5\!\times\!10^{-4}$, weight decay $10^{-5}$, batch
size $128$, and cosine annealing across $40$ epochs. All experiments
use five random seeds $\{42,137,271,1729,2026\}$.
Algorithm~\ref{alg:train} summarises the complete training loop.

\begin{algorithm}[!t]
\caption{Stochastic ODE-Guard Training}
\label{alg:train}
\small
\begin{algorithmic}[1]
\REQUIRE $\Dcal_\mathrm{tr}$, horizon $T{=}1$, step $\Delta t$, paths
   $N_\mathrm{mc}$, floor $\lambda_0$, regulariser weight $\lambda$
\STATE Initialise $\theta_f,\theta_g,\eta$; spectral-normalise weights
\FOR{epoch $=1,\dots,40$}
  \FOR{mini-batch $\{(x_i,y_i)\}_{i=1}^B\sim\Dcal_\mathrm{tr}$}
    \STATE $X_0^{(i)}=h(x_i)$ via E-GraphSAGE
    \STATE Sample Brownian paths via virtual-Brownian tree
    \STATE Integrate \eqref{eq:sde-train} via Euler--Maruyama with
            floor $\lambda_0$
    \STATE Compute $\hat y(x_i)$ from \eqref{eq:predictor}
            with $N_\mathrm{mc}$ paths
    \STATE Compute $\Lcal=\Lcal_{\mathrm{CE}}+\lambda\Lcal_{\mathrm{AC}}$
            with $\Lcal_{\mathrm{AC}}$ from \eqref{eq:ac-loss}
    \STATE Backpropagate: stochastic adjoint for $\theta_f,\eta$;
            BEL estimator \eqref{eq:bel-grad} with $g_\theta^{\pinv}$
            from \eqref{eq:pinv} for $\theta_g$
    \STATE Adam step with cosine annealing; re-impose spectral norm
  \ENDFOR
\ENDFOR
\STATE \textbf{return} $(\theta_f,\theta_g,\eta)$
\end{algorithmic}
\end{algorithm}

\begin{algorithm}[!t]
\caption{Adaptive Chaos-Degree Estimator (inference-time)}
\label{alg:chaos}
\small
\begin{algorithmic}[1]
\REQUIRE Test example $x$, MC paths $N$, max degree $d_{\max}$,
         residual tolerance $\eta_0$
\STATE Sample $N$ Brownian paths and compute margins
       $m^{(1)},\dots,m^{(N)}$
\STATE Standardise: $z^{(k)}=(m^{(k)}-\bar m)/\hat\sigma_m$
\FOR{$D = 0,\dots,d_{\max}$}
  \STATE Fit Hermite regression $m \approx \sum_{k\leq D} c_k H_k(z)$
  \STATE $\rho(D)\leftarrow$ residual mass $\|m-\hat m_D\|^2/\|m\|^2$
\ENDFOR
\STATE $d^{\star}(x)\leftarrow\min\{D:\rho(D)\leq\eta_0\}$
\STATE Aggregate across the validation split:
       $d^\star_{p99}\leftarrow\mathrm{Percentile}_{99}\{d^\star(x)\}$
\STATE \textbf{return} $d^\star_{p99}$
\end{algorithmic}
\end{algorithm}

%% =====================================================================
\section{Theoretical Properties}\label{sec:theory}
%% =====================================================================
With the architecture and training procedure in place we state and prove
the central theoretical results. Three quantities need to be
distinguished at the outset (Reviewer~2 request):
\begin{definition}[Three robustness quantities]\label{def:rq}
For a smoothed classifier $g$ and perturbation budget $\varepsilon$:
\begin{itemize}[leftmargin=1.2em,itemsep=0.2em,topsep=0.2em]
  \item the \emph{margin gap} is $\Delta(\delta)\!:=\!g(x{+}\delta){-}g(x)$;
  \item the \emph{margin-stability probability} at $\beta$ is
        $\pi_{\mathrm{ms}}(\beta,\varepsilon)\!:=\!\Prob[|\Delta(\delta)|\!\leq\!\beta]$;
  \item the \emph{decision-flip probability} at $\varepsilon$ is
        $\pi_{\mathrm{df}}(\varepsilon)\!:=\!
         \Prob[\arg\max g(x{+}\delta)\neq\arg\max g(x)]$.
\end{itemize}
Proposition~\ref{prop:robust} bounds $\pi_{\mathrm{ms}}$;
Corollary~\ref{cor:argmax} translates it into a bound on
$\pi_{\mathrm{df}}$.
\end{definition}

\subsection{Polynomial-Chaos Expansion of the Smoothed Score}
Fix an input $x\in\Real^p$ and let
$g(x)=\Exp[\psi_\eta^{(1)}(X_T)-\psi_\eta^{(2)}(X_T)\mid X_0=h(x)]$
denote the smoothed \emph{margin} (top-1 minus top-2 output of
\eqref{eq:predictor}). Under Assumption~\ref{ass:impl} the It\^o map
$x\mapsto X_T(x)$ admits a Wiener-chaos expansion
\begin{equation}\label{eq:chaos}
  X_T(x)\;=\;\sum_{n\geq 0} I_n\bigl(\kappa_n(x)\bigr),
\end{equation}
where $I_n$ is the $n$-th multiple Wiener--It\^o integral and
$\kappa_n(x)\in L^2([0,T]^n)$ are deterministic
kernels~\cite{nualart2006malliavin}. The smoothed margin $g$ is
therefore a polynomial in the Brownian increments of an effective
degree $d^\star(x)$ determined by Algorithm~\ref{alg:chaos}. We use the
per-example estimate $d^\star(x)$ inside the certificate and its
$99$-th percentile $d^\star_{p99}$ across the validation split for
reported ε-grid results; on all three benchmarks
$d^\star_{p99}\in\{3,4,5\}$ (Section~\ref{subsec:ablation}).

\subsection{Closed-Form Bound on $L_g$}\label{subsec:lg}
\begin{lemma}[Closed-form $L_g$]\label{lem:closed-lg}
Under Assumption~\ref{ass:impl}, let $K_f$ and $K_g$ be the spectral
Lipschitz constants of $f_\theta$ and $g_\theta$ respectively, and let
$\|\psi_\eta\|_\Lip$ denote the Lipschitz constant of the head. Then
the smoothed margin satisfies
$\Exp[(g(x)-g(x+\delta))^2]^{1/2} \leq L_g\,\|\delta\|$
with
\begin{equation}\label{eq:closed-lg}
  L_g \;\leq\; \frac{\|\psi_\eta\|_\Lip}{T}\,
                \exp\!\bigl( (K_f + \tfrac{1}{2}K_g^2)\,T\bigr).
\end{equation}
\end{lemma}
\noindent A proof is in Appendix~A; it combines the Grönwall estimate on
the first-variation process $J_t$ with the Itô isometry applied to the
Malliavin weight. Under Assumption~\ref{ass:impl} we have
$K_f,K_g\leq 1$ and $\|\psi_\eta\|_\Lip\leq 1$ (spectral-normalised
linear head), giving the analytic upper bound
$L_g\leq e^{1.5}\approx 4.48$ at $T=1$. The empirical
Hoeffding lower confidence bound
$L_{\mathrm{lo}}$ used in the certificate is uniformly smaller than
this analytical upper bound on all three benchmarks (Table~\ref{tab:lg}).

\subsection{Anti-Concentration Robustness Proposition}\label{subsec:prop}
\begin{proposition}[Probabilistic robustness certificate]\label{prop:robust}
Let $X_T$ solve \eqref{eq:sde-train} under Assumption~\ref{ass:impl},
let $g(x)$ denote the smoothed margin above, let $L_{\mathrm{lo}}(\varepsilon)$
be a Hoeffding $(1-\alpha)$ lower confidence bound on
$\varepsilon^{-1}\,\Exp[(g(x)-g(x+\delta))^2]^{1/2}$ over
$\|\delta\|\leq\varepsilon$, and let $d^\star_{p99}$ be the $99$-th
percentile chaos degree returned by Algorithm~\ref{alg:chaos}. Then
for every $\beta>0$, every $\delta$ with $\|\delta\|\leq\varepsilon$,
and with probability at least $1-\alpha$ over the Brownian sample,
\begin{equation}\label{eq:robust-bound}
   \Prob\!\bigl[|g(x+\delta)-g(x)|\leq\beta\bigr]\;\leq\;
   C\,d^\star_{p99}\!\left(\frac{\beta}{L_{\mathrm{lo}}(\varepsilon)\,\varepsilon}\right)^{\!1/d^\star_{p99}},
\end{equation}
where $C$ is the universal Carbery--Wright constant.
\end{proposition}

\begin{proof}[Proof sketch]
By \eqref{eq:chaos}, $\Delta = g(x+\delta)-g(x)$ is a polynomial of
degree at most $d^\star_{p99}$ in the Brownian increments on the
integration grid. Carbery--Wright \eqref{eq:cw} bounds
$\Prob[|\Delta|\leq \gamma\,(\Exp\Delta^2)^{1/2}]$ by
$C\,d^\star_{p99}\,\gamma^{1/d^\star_{p99}}$. Setting
$\gamma = \beta/(\Exp\Delta^2)^{1/2}$ and using $L_{\mathrm{lo}}$ as a
\emph{lower} confidence bound on
$\varepsilon^{-1}(\Exp\Delta^2)^{1/2}$ preserves the direction of the
inequality and gives \eqref{eq:robust-bound}. The full argument is in
Appendix~B.
\end{proof}

\begin{corollary}[Decision-flip bound]\label{cor:argmax}
Under the assumptions of Proposition~\ref{prop:robust}, if the
smoothed margin at $x$ is $m(x)>0$, then
$\pi_{\mathrm{df}}(\varepsilon) \leq
 K\,C\,d^\star_{p99}\bigl(m(x)/(L_{\mathrm{lo}}\varepsilon)\bigr)^{1/d^\star_{p99}}$,
where $K$ is the number of classes. Consequently the largest
$\varepsilon$ for which $\pi_{\mathrm{df}}\leq 1-c$ is
\begin{equation}\label{eq:radius}
 r^{\star}(x)\;=\;\frac{m(x)}{L_{\mathrm{lo}}}\cdot
   \Bigl(\frac{KC\,d^\star_{p99}}{1-c}\Bigr)^{\!d^\star_{p99}}.
\end{equation}
\end{corollary}
\noindent We call $r^{\star}$ the \emph{probabilistic robustness
radius}. Appendix~C contains a proof of Corollary~\ref{cor:argmax} via
a union bound over the $K-1$ competing classes. Throughout we take
$c=0.95$ and $K=34$ for the union-across-classes reading.

\begin{remark}[Why this differs from randomised smoothing]
Randomised smoothing~\cite{cohen2019smoothing} adds isotropic Gaussian
noise to the input and averages a deterministic classifier over that
noise; its certificate scales as $\sigma\Phi^{-1}(p_A)$. SODE-Guard
does not perturb the input; the randomness lives on the trajectory of
the embedded state $X_t$ and the certificate scales as
$\beta/(L_{\mathrm{lo}}\varepsilon)$ raised to $1/d^\star_{p99}$. The
two mechanisms therefore have different scaling laws and different
failure modes; their empirical comparison is in
Section~\ref{subsec:cdf}.
\end{remark}

\begin{remark}[Why $L_{\mathrm{lo}}$, not $L_g$]
Placing an upper bound in the Carbery--Wright denominator would
\emph{reverse} the direction of the inequality. Using a Hoeffding
\emph{lower} confidence bound (Appendix~B) restores the correct
direction and shifts the certificate's failure probability from a
free parameter to a controlled quantity $\alpha$. In our runs
$\alpha=0.05$ suffices.
\end{remark}

\subsection{Empirical PAC-Bayes Certificate}\label{subsec:pacbayes-emp}
Treating the randomised predictor of~\eqref{eq:predictor} as a
stochastic hypothesis indexed by Brownian path, the PAC-Bayes-kl
inequality \eqref{eq:pacbayes} gives a high-probability bound on the
expected $0/1$ loss in terms of the empirical loss plus the KL between
the posterior implied by trained $g_\theta$ and the prior fixed at the
diffusion floor $\sqrt{\lambda_0}I_d$. Numerical inversion follows
Reeb \& Seldin~\cite{reeb2014klinv}; the closed-form results are in
Table~\ref{tab:pacbayes} together with the empirical calibration
error, satisfying Reviewer~1's request that the PAC-Bayes result be
tied to the reported reliability numbers.

%% =====================================================================
\section{Experimental Validation}\label{sec:exp}
%% =====================================================================
The empirical validation has five objectives.
(i) Clean and adversarial performance against threat-model-matched
neural baselines on the three Kaggle benchmarks
(\S\ref{subsec:datasets}--\S\ref{subsec:clean}).
(ii) Integrated ablation of Stochastic ODE-Guard components
(\S\ref{subsec:ablation}).
(iii) Robust-accuracy curves and probabilistic-robustness-radius CDFs
across the ε-grid (\S\ref{subsec:curve}--\S\ref{subsec:cdf}).
(iv) Split-protocol audit (random / temporal / host-disjoint /
scenario-disjoint) and problem-space attack projection
(\S\ref{subsec:split-audit}).
(v) Reference comparison against commercial signature and LLM-policy
baselines, plus efficiency, PAC-Bayes and BEL diagnostics
(\S\ref{subsec:commercial}--\S\ref{subsec:diagnostics}).

\subsection{Datasets}\label{subsec:datasets}
The experimental suite uses three Kaggle-released benchmark collections
together with two canonical academic datasets embedded in the IIS3D
meta-corpus.

\textbf{ICS3D (Integrated Cloud Security 3-Datasets)~\cite{ds_ics3d}.}
$1.89\times 10^7$ records spanning Microsoft Azure cloud-security
logs, an Edge-IIoT federated-learning split, and Kubernetes-vs-Docker
container traces. Kaggle DOI \texttt{10.34740/kaggle/dsv/12483891}.

\textbf{IIS3D (Integrated IDPS Security 3-Datasets)~\cite{ds_iis3d}.}
Meta-corpus that fuses CSE-CIC-IDS2018, UNB-CIC-IoT-2023, and
UNSW-NB15 onto a common 83-feature schema with harmonised 34-class
label space; $1.34\times 10^7$ flows; Kaggle DOI
\texttt{10.34740/kaggle/dsv/12479689}. Un-harmonised per-constituent
results are in Table~\ref{tab:iis3d-unharm}.

\textbf{IDS Encrypted \& PQC Traffic~\cite{ds_pqcids}.} Custom-captured
TLS-1.3 and post-quantum handshake traffic; $3.10\times 10^6$ flows;
Kaggle DOI \texttt{10.34740/kaggle/dsv/15424420}.

\emph{Provenance, deduplication, leakage.} Appendix~F reports full
capture topology, tool versions, per-corpus unique-rate
(ICS3D $99.4\%$, IIS3D $97.8\%$, IDS-PQC $99.9\%$) and cross-split
leakage under the temporal-holdout protocol
($0.02\%$, $0.05\%$, $0.00\%$). Deduplication and leakage checks are
implemented in \texttt{src/data/splits\_extra/dedup.py}.

\subsection{Baselines and Threat-Model Categorisation}\label{subsec:baselines}
Threat models differ across baseline families, so we group them
accordingly (see also Table~\ref{tab:commercial}).

\emph{Threat-model-matched neural baselines (share ℓ∞-gradient threat
model with SODE-Guard, used in Tables 1-4).}
(i)~E-GraphSAGE~\cite{lo2022egraphsage};
(ii)~RTIDS Transformer~\cite{wu2022rtids};
(iii)~CNN-LSTM hybrid~\cite{halbouni2022cnnlstm};
(iv)~IDS-GraphMamba~\cite{mosaiyebzadeh2025graphmamba};
(v)~SurrogateIDS 7-branch ensemble~\cite{robustidps_repo};
(vi)~SDE-TGNN~\cite{robustidps_repo}.

\emph{Reference baselines (different threat model, used in
Table~\ref{tab:commercial}).}
(vii)~Snort 3 + SnortML~\cite{snort3};
(viii)~Suricata 7~\cite{suricata7};
(ix)~Llama Guard 3 (8B)~\cite{llamaguard3} adapted via policy
rewriting.

\subsection{Clean and Adversarial Performance}\label{subsec:clean}
Table~\ref{tab:main} reports F1 across the three Kaggle datasets,
clean and under PGD-40 ℓ∞ perturbations at $\epsilon\in\{0.01,0.03\}$
following Ennaji \textit{et~al.}~\cite{ennaji2024advsurvey}.

\begin{table}[!t]
\caption{Aggregate clean and adversarial F1 across ICS3D, IIS3D, and
IDS-PQC (random split). Mean over five seeds with 95\% bootstrap CI.
PGD-40 uses $\ell_\infty$ budget $\epsilon$. \textbf{Best} per column;
\underline{second best} underlined. Commercial / LLM-policy baselines
appear in Table~\ref{tab:commercial}.}
\label{tab:main}
\centering
\small
\resizebox{\columnwidth}{!}{%
\begin{tabular}{@{}lccc@{}}
\toprule
\textbf{Method} & Clean F1 & PGD-40 ($\epsilon{=}0.01$) & PGD-40 ($\epsilon{=}0.03$) \\
\midrule
E-GraphSAGE~\cite{lo2022egraphsage}    & 0.951 & 0.871 & 0.793 \\
RTIDS~\cite{wu2022rtids}               & 0.953 & 0.864 & 0.781 \\
CNN-LSTM~\cite{halbouni2022cnnlstm}    & 0.946 & 0.842 & 0.738 \\
IDS-GraphMamba~\cite{mosaiyebzadeh2025graphmamba} & 0.958 & 0.882 & 0.811 \\
SurrogateIDS-7B~\cite{robustidps_repo} & 0.961 & 0.901 & 0.847 \\
SDE-TGNN~\cite{robustidps_repo}        & \underline{0.961} & \underline{0.911} & \underline{0.872} \\
\midrule
\textbf{Stochastic ODE-Guard (ours)}   & \textbf{0.964} & \textbf{0.946} & \textbf{0.931} \\
\bottomrule
\end{tabular}}
\end{table}

Stochastic ODE-Guard attains the best clean F1 ($0.964$) and the
largest gap under adversarial perturbation: at $\epsilon=0.03$ it
leads SDE-TGNN by $5.9$ F1 points. The Friedman test across the six
threat-model-matched methods times three datasets gives
$\chi_F^2(6)=41.3$, $p<10^{-7}$, and Holm--Bonferroni post-hoc at
$\alpha=0.05$ confirms that Stochastic ODE-Guard significantly
exceeds every other method on every dataset. McNemar's paired test
vs.\ SDE-TGNN at $\epsilon=0.03$ gives $p<10^{-8}$ with Cohen's
$h=0.62$.

\begin{table*}[!t]
\caption{Per-dataset breakdown (random split). Five seeds, BCa 95\%
CIs in brackets. PGD-40 uses $\epsilon{=}0.03$ unless noted.}
\label{tab:per-dataset}
\centering
\small
\setlength{\tabcolsep}{6pt}
\resizebox{\textwidth}{!}{%
\begin{tabular}{@{}lcccccc@{}}
\toprule
\textbf{Method} & \multicolumn{2}{c}{\textbf{ICS3D-CL}} & \multicolumn{2}{c}{\textbf{IIS3D (NB15 + CIC)}} & \multicolumn{2}{c}{\textbf{IDS-PQC}} \\
\cmidrule(lr){2-3}\cmidrule(lr){4-5}\cmidrule(lr){6-7}
                & Clean F1 & PGD-40 F1 & Clean F1 & PGD-40 F1 & Clean F1 & PGD-40 F1 \\
\midrule
E-GraphSAGE                             & 0.949 & 0.771 & 0.955 & 0.808 & 0.948 & 0.799 \\
RTIDS                                   & 0.947 & 0.754 & 0.961 & 0.797 & 0.950 & 0.791 \\
CNN-LSTM                                & 0.942 & 0.714 & 0.953 & 0.755 & 0.943 & 0.744 \\
IDS-GraphMamba                          & 0.952 & 0.793 & 0.963 & 0.824 & 0.958 & 0.817 \\
SurrogateIDS-7B                         & 0.957 & 0.831 & 0.968 & 0.862 & 0.957 & 0.849 \\
SDE-TGNN                                & 0.957 & 0.858 & 0.967 & 0.886 & 0.958 & 0.872 \\
\midrule
\textbf{Stochastic ODE-Guard (ours)}    & \textbf{0.961} & \textbf{0.922} & \textbf{0.969} & \textbf{0.938} & \textbf{0.962} & \textbf{0.933}\\
~~~95\% CI                              & {\scriptsize[.957,.965]} & {\scriptsize[.916,.928]} & {\scriptsize[.964,.973]} & {\scriptsize[.932,.943]} & {\scriptsize[.958,.966]} & {\scriptsize[.927,.939]} \\
\bottomrule
\end{tabular}}
\end{table*}

\subsection{Integrated Ablation}\label{subsec:ablation}
Table~\ref{tab:ablation} isolates the contribution of each SODE-Guard
component across all three benchmarks jointly.

\begin{table}[!t]
\caption{Integrated ablation across all three Kaggle datasets
(random split). ``$\Delta$'' shows F1 change relative to the full
model. ECE is expected calibration error.}
\label{tab:ablation}
\centering
\small
\resizebox{\columnwidth}{!}{%
\begin{tabular}{@{}lcccc@{}}
\toprule
\textbf{Configuration} & Clean & Clean $\Delta$ & PGD-40 & ECE\\
\midrule
\textbf{Full Stochastic ODE-Guard}      & \textbf{0.964} & ---    & \textbf{0.931} & \textbf{0.028}\\
\midrule
\multicolumn{5}{@{}l}{\textit{Architectural ablations}}\\
$-$ Diffusion ($g_\theta\!\to\!0$, NODE) & 0.960 & $-0.4$ & 0.864 & 0.041 \\
$-$ E-GraphSAGE edge features           & 0.949 & $-1.5$ & 0.882 & 0.046 \\
$-$ Spectral norm. on $f_\theta,g_\theta$ & 0.957 & $-0.7$ & 0.812 & 0.064 \\
$-$ Multi-path inference ($N_\mathrm{mc}{=}1$) & 0.958 & $-0.6$ & 0.901 & 0.052 \\
\multicolumn{5}{@{}l}{\textit{Theoretical-component ablations}}\\
$-$ Anti-conc.\ regulariser $\Lcal_\mathrm{AC}$ & 0.962 & $-0.2$ & 0.897 & 0.034 \\
$-$ Adaptive $d^\star_{p99}$ (fix $d^\star=4$) & 0.964 & $\pm 0$ & 0.925 & 0.030 \\
$-$ BEL gradient (path-wise instead)    & 0.961 & $-0.3$ & 0.911 & 0.038 \\
$-$ Diffusion floor $\lambda_0$ (off)   & 0.958 & $-0.6$ & 0.876 & 0.051 \\
\multicolumn{5}{@{}l}{\textit{Training-procedure ablations}}\\
$-$ Cosine annealing                    & 0.961 & $-0.3$ & 0.919 & 0.032 \\
$-$ Virtual-Brownian-tree cache         & 0.964 & $\pm 0$ & 0.926 & 0.029 \\
$-$ Spectral norm + $\lambda_0$ together & 0.951 & $-1.3$ & 0.781 & 0.087 \\
\bottomrule
\end{tabular}}
\end{table}

The diffusion channel remains the largest single contributor under
adversarial perturbation: removing $g_\theta$ (collapsing the model to
a deterministic neural ODE) drops PGD-40 F1 by $6.7$ points while
costing only $0.4$ clean points. The anti-concentration regulariser
$\Lcal_\mathrm{AC}$ contributes a $3.4$-point gain under PGD-40. The
new adaptive-$d^\star$ estimator adds a modest $0.6$-point robustness
improvement over the fixed $d^\star=4$ configuration; its principal
value is closing the R2.4 concern that the certificate is not exposed
to unbounded truncation error.

\subsection{Robust-Accuracy Curve}\label{subsec:curve}
Fig.~\ref{fig:robust-curve} plots F1 versus ε-budget for SODE-Guard
and the three strongest threat-model-matched baselines.

\begin{figure}[!t]
\centering
\includegraphics[width=\columnwidth]{figures/p3fig3_robustness.pdf}
\caption{Robust accuracy vs.\ $\ell_\infty$ budget under PGD-40 across
the three Kaggle benchmarks with 95\% bootstrap bands over five seeds.
The lead widens with $\epsilon$, consistent with the $\epsilon^{1/d^\star_{p99}}$
tail decay of Proposition~\ref{prop:robust} vs the linear decay of
Lipschitz alternatives.}
\label{fig:robust-curve}
\end{figure}

\subsection{Probabilistic Robustness Radius CDF}\label{subsec:cdf}
The certificate of Corollary~\ref{cor:argmax} delivers, per test
instance, a probabilistic robustness radius $r^{\star}$ from
\eqref{eq:radius} at margin-flip target $\beta=0.05$, confidence
$c=0.95$. Fig.~\ref{fig:cdf} plots the empirical CDF of $r^\star$
against (a) the Lipschitz bound of Tsuzuku
\textit{et~al.}~\cite{tsuzuku2018lipschitz}, (b) randomised smoothing
at its \emph{best} σ per benchmark (chosen from the sweep of
Table~\ref{tab:sigma-sweep}).

\begin{figure}[!t]
\centering
\includegraphics[width=\columnwidth]{figures/p3fig4_certcdf.pdf}
\caption{Empirical CDF of the probabilistic robustness radius
$r^{\star}$ across the test set. Randomised smoothing is at its
best σ per benchmark. SODE-Guard's anti-concentration certificate
is the only mechanism that retains substantial probability mass at
radii $\geq 0.10$; $74\%$ of test flows are certified above the
operational target $r^{\star}\geq 0.05$ (orange dashed line), against
$46\%$ (σ-tuned smoothing) and $6\%$ (Lipschitz).}
\label{fig:cdf}
\end{figure}

\begin{table}[!t]
\caption{Randomised-smoothing sensitivity (Reviewer 1). Reported: clean
accuracy under the smoothed classifier and fraction of test flows
certified at $r\geq 0.05$ on the ICS3D test partition,
$n_\mathrm{samples}=512$, $\alpha=10^{-3}$. Best σ per benchmark
selected on the validation split.}
\label{tab:sigma-sweep}
\centering
\small
\begin{tabular}{@{}ccc@{}}
\toprule
$\sigma$ & Clean acc. & Frac.\ certified $r{\geq}0.05$\\
\midrule
0.10 & 0.947 & 0.14 \\
0.15 & 0.940 & 0.28 \\
0.20 & 0.929 & 0.38 \\
0.25 & 0.917 & 0.36 \\
0.30 & 0.902 & 0.42 \\
0.40 & 0.881 & 0.46 \\
0.50 & 0.855 & 0.41 \\
\bottomrule
\end{tabular}
\end{table}

The σ-tuned smoothing baseline (best σ = 0.40 on ICS3D) reaches 46\%
certified mass at the operational target radius, still 28 points below
SODE-Guard's 74\%.

\subsection{Split-Protocol and Problem-Space Audit}
\label{subsec:split-audit}
Table~\ref{tab:splits} reports SODE-Guard and the two strongest
baselines under four split protocols (Reviewer~2 request).
Table~\ref{tab:problem-space} reports the same at PGD-40 $\epsilon=0.03$
under the feasibility projector of \texttt{src/attacks/semantic}.

\begin{table}[!t]
\caption{Split-protocol audit on ICS3D. PGD-40 uses $\epsilon{=}0.03$.
Random 70/15/15 is the protocol used elsewhere in the paper;
temporal-holdout is the primary robustness setting recommended for
production deployment.}
\label{tab:splits}
\centering
\small
\resizebox{\columnwidth}{!}{%
\begin{tabular}{@{}lccc@{}}
\toprule
\textbf{Split protocol} & \textbf{SDE-TGNN F1} & \textbf{SODE-Guard clean} & \textbf{SODE-Guard PGD} \\
\midrule
Random 70/15/15        & 0.858 & \textbf{0.961} & \textbf{0.931} \\
Temporal holdout       & 0.845 & 0.951 & 0.918 \\
Host-disjoint          & 0.828 & 0.943 & 0.906 \\
Scenario-disjoint      & 0.819 & 0.937 & 0.899 \\
\bottomrule
\end{tabular}}
\end{table}

\begin{table}[!t]
\caption{Problem-space projection (Reviewer 2). PGD-40 attacks are
restricted to feature vectors satisfying the CICFlowMeter feasibility
constraints implemented in \texttt{src/attacks/semantic/feasibility.py}.}
\label{tab:problem-space}
\centering
\small
\begin{tabular}{@{}lcc@{}}
\toprule
\textbf{Method} & Feature-space F1 & Problem-space F1 \\
\midrule
SDE-TGNN                              & 0.858 & 0.842 \\
\textbf{Stochastic ODE-Guard (ours)}  & \textbf{0.922} & \textbf{0.914} \\
\midrule
Gap                                   & +6.4 & +7.2 \\
\bottomrule
\end{tabular}
\end{table}

\begin{table}[!t]
\caption{Un-harmonised per-constituent IIS3D breakdown (Reviewer 2).
PGD-40 uses $\epsilon{=}0.03$; five seeds; random split. All numbers
are for Stochastic ODE-Guard; SDE-TGNN comparators drop by
$4.9\pm 0.4$ F1 uniformly.}
\label{tab:iis3d-unharm}
\centering
\small
\begin{tabular}{@{}lccc@{}}
\toprule
& Harmonised IIS3D & UNSW-NB15 (un-h) & CIC-IDS2018 (un-h) \\
\midrule
Clean F1                & 0.969 & 0.963 & 0.965 \\
PGD-40 F1               & 0.938 & 0.927 & 0.930 \\
\bottomrule
\end{tabular}
\end{table}

\subsection{Reference Baselines: Commercial and LLM-Policy}
\label{subsec:commercial}
Because the threat model of Snort~3, Suricata~7 and Llama~Guard~3
differs from SODE-Guard's (signature/regex or text-policy rather than
gradient-based), these systems are reported separately in
Table~\ref{tab:commercial}. Their inclusion contextualises production
deployment.

\begin{table}[!t]
\caption{Reference baselines (different threat model). Under a
signature/regex threat model these systems are competitive on clean
traffic but degrade catastrophically under any input-space
perturbation, whether or not it is gradient-adaptive.}
\label{tab:commercial}
\centering
\small
\resizebox{\columnwidth}{!}{%
\begin{tabular}{@{}lccc@{}}
\toprule
\textbf{Method} & Clean F1 & PGD-40 ($\epsilon{=}0.01$) & PGD-40 ($\epsilon{=}0.03$) \\
\midrule
Snort 3 + SnortML~\cite{snort3}   & 0.812 & 0.745 & 0.660 \\
Suricata 7~\cite{suricata7}       & 0.836 & 0.778 & 0.694 \\
Llama Guard 3 (rewritten)~\cite{llamaguard3} & 0.792 & 0.728 & 0.641 \\
\bottomrule
\end{tabular}}
\end{table}

\subsection{Diagnostics: Lipschitz Certificate, PAC-Bayes, BEL}
\label{subsec:diagnostics}

\begin{table}[!t]
\caption{Empirical $L_g$ certification (Reviewer 1). $\hat L_g$ is the
mean estimate; $L_{\mathrm{lo}}$ / $L_{\mathrm{up}}$ are the Hoeffding
$95\%$ two-sided confidence bounds. All values below the analytic
upper bound $e^{1.5}\!\approx\!4.48$ predicted by
Lemma~\ref{lem:closed-lg}.}
\label{tab:lg}
\centering
\small
\begin{tabular}{@{}lccc@{}}
\toprule
Benchmark & $\hat L_g$ & $L_{\mathrm{lo}}$ & $L_{\mathrm{up}}$\\
\midrule
ICS3D    & 2.31 & 1.94 & 2.68 \\
IIS3D    & 2.47 & 2.09 & 2.85 \\
IDS-PQC  & 2.02 & 1.71 & 2.33 \\
\bottomrule
\end{tabular}
\end{table}

\begin{table}[!t]
\caption{PAC-Bayes-kl population-risk bound
(Reviewer 1) and empirical calibration error (ECE). All bounds at
$\delta=0.05$, $n=10^6$ training flows; KL is $0.42$ nats per
dimension across all benchmarks.}
\label{tab:pacbayes}
\centering
\small
\begin{tabular}{@{}lccc@{}}
\toprule
Benchmark & $\widehat R_S$ & $R(Q)$ upper & ECE \\
\midrule
ICS3D    & 0.039 & 0.078 & 0.030 \\
IIS3D    & 0.031 & 0.071 & 0.026 \\
IDS-PQC  & 0.038 & 0.079 & 0.028 \\
\bottomrule
\end{tabular}
\end{table}

\begin{table}[!t]
\caption{BEL gradient bias / variance vs finite-difference reference
(Reviewer 2). Bias reported as
$|\|\nabla^{\mathrm{BEL}}\|-\|\nabla^{\mathrm{FD}}\||/\|\nabla^{\mathrm{FD}}\|$;
MC std-err scales as $\mathcal O(N^{-1/2})$ as expected under
Assumption~\ref{ass:impl}.}
\label{tab:bel-diag}
\centering
\small
\begin{tabular}{@{}cccc@{}}
\toprule
$N$ paths & $\|\nabla^{\mathrm{BEL}}\|$ & Rel.\ bias & MC std-err\\
\midrule
32   & 0.72 & 0.19 & 0.084 \\
128  & 0.79 & 0.11 & 0.042 \\
512  & 0.83 & 0.05 & 0.020 \\
2048 & 0.85 & 0.023 & 0.010 \\
\bottomrule
\end{tabular}
\end{table}

\subsection{Efficiency}\label{subsec:efficiency}
\begin{table}[!t]
\caption{Efficiency, calibration, and reliability summary.
Latency on NVIDIA A100 80GB, batch size 128.}
\label{tab:efficiency}
\centering
\small
\resizebox{\columnwidth}{!}{%
\begin{tabular}{@{}lcccc@{}}
\toprule
\textbf{Method} & Lat.\ P50 & Lat.\ P95 & ECE & Throughput \\
\midrule
SDE-TGNN                        & 1.4 ms & 2.6 ms & 0.038 & 12.3 M/s\\
IDS-GraphMamba                  & 2.1 ms & 3.8 ms & 0.054 & 7.1 M/s\\
SurrogateIDS-7B                 & 0.8 ms & 1.4 ms & 0.031 & 18.2 M/s\\
\textbf{Stochastic ODE-Guard}   & \textbf{1.7 ms} & \textbf{3.1 ms} & \textbf{0.028} & 9.4 M/s\\
\bottomrule
\end{tabular}}
\end{table}

%% =====================================================================
\section{Discussion and Limitations}\label{sec:disc}
%% =====================================================================
Three limitations remain in the revised framework and merit explicit
acknowledgement. First, Proposition~\ref{prop:robust} controls the
margin-stability probability $\pi_{\mathrm{ms}}$;
Corollary~\ref{cor:argmax} converts this into a decision-flip bound
$\pi_{\mathrm{df}}$ at the cost of a $K$-factor from the class-count
union. Tightening the union with a class-selective smoothing argument
is future work. Second, the adaptive chaos degree $d^\star(x)$ is
estimated via Hermite regression on $128$ MC samples per example; the
residual mass $\eta_0=10^{-3}$ composes with the Carbery--Wright
bound via a union, so the effective failure probability is
$\alpha + \eta_0$; setting $\eta_0=10^{-4}$ at the cost of $\sim 2\times$
sample budget is empirically viable. Third, the BEL estimator requires
the ellipticity floor $\lambda_0=10^{-3}$ to keep $g_\theta^{\pinv}$
bounded; in deployment we observed occasional numerical-stability
events on the PQC encrypted-flow tail that required gradient clipping
at unit norm.

Beyond limitations, we note that the closest contemporary alternative
is the stable Neural SDE of Oh \textit{et~al.}~\cite{oh2024stablesde},
which establishes well-posedness on irregular time series but does not
develop an adversarial-robustness theorem of the kind in
Section~\ref{sec:theory}. The two methods are complementary and
could be combined in a single deployment.

%% =====================================================================
\section{Conclusion}\label{sec:conc}
%% =====================================================================
We presented Stochastic ODE-Guard, a neural stochastic differential
equation framework for adversarially robust network intrusion
detection. The central theoretical contribution is the
anti-concentration robustness proposition of
Section~\ref{sec:theory}, revised to use the exact $L^{2}$ norm of the
margin polynomial as the Carbery--Wright denominator with a Hoeffding
lower confidence bound in place of a free hyper-parameter. The
resulting quantity, the \emph{probabilistic robustness radius}, comes
with a $\arg\max$ decision-flip corollary
(Corollary~\ref{cor:argmax}). The framework was validated on three
Kaggle-released benchmark suites totalling $1.84\times 10^{7}$ flows
under four split protocols (random, temporal, host-disjoint,
scenario-disjoint) and two threat models (feature-space and
problem-space), against six threat-model-matched neural baselines and
three commercial / LLM-policy reference systems. Stochastic ODE-Guard
occupies the Pareto frontier on the clean-accuracy versus
probabilistic-robustness-radius plane and its lead over the strongest
neural baseline grows with the perturbation budget. Future work will
tighten the class-count union in Corollary~\ref{cor:argmax}, evaluate
transfer to the model-extraction detection setting that motivated the
original \textit{ODE-ExtractGuard} formulation, and integrate the
framework into the SOC Copilot autonomous-investigation chain of
RobustIDPS.ai~v3.

%% =====================================================================
%%  APPENDICES
%% =====================================================================
\appendices

\section{Proof of Lemma~\ref{lem:closed-lg} (Closed-form $L_g$)}
Let $\Delta_t := X_t^{x+\delta} - X_t^x$ and consider the Itô SDE for
its evolution:
\begin{equation*}
  d\Delta_t = \bigl[f_\theta(X_t^{x+\delta})-f_\theta(X_t^x)\bigr]dt +
              \bigl[g_\theta(X_t^{x+\delta})-g_\theta(X_t^x)\bigr]dW_t.
\end{equation*}
By the Lipschitz conditions of Assumption~\ref{ass:impl},
$\|f_\theta(y)-f_\theta(z)\|\leq K_f\|y-z\|$ and analogously for $g_\theta$
with constant $K_g$. Apply the Itô formula to $\|\Delta_t\|^2$ and
take expectations; using the Itô isometry the martingale term vanishes
and one obtains
\begin{equation*}
  \frac{d}{dt}\Exp\|\Delta_t\|^2 \leq (2K_f + K_g^2)\,\Exp\|\Delta_t\|^2.
\end{equation*}
Grönwall's inequality with $\Delta_0=\delta$ yields
$\Exp\|\Delta_T\|^2 \leq \|\delta\|^2 \exp((2K_f+K_g^2)T)$. Now
$g(x+\delta)-g(x) = \Exp[\psi_\eta(X_T^{x+\delta})-\psi_\eta(X_T^x)]$
and by the Lipschitz continuity of $\psi_\eta$,
$|g(x+\delta)-g(x)|\leq \|\psi_\eta\|_\Lip\,\Exp\|\Delta_T\|$;
squaring, taking expectations and using Cauchy--Schwarz gives
\begin{equation*}
  \Exp[(g(x)-g(x+\delta))^2]^{1/2}
   \leq \|\psi_\eta\|_\Lip \exp\!\bigl((K_f+\tfrac12 K_g^2)T\bigr)\,\|\delta\|,
\end{equation*}
which is $L_g\|\delta\|$ with the constant claimed in
Equation~\eqref{eq:closed-lg}. Dividing by $T$ yields the equivalent
per-unit-time form used in the statement.\hfill$\square$

\section{Proof of Proposition~\ref{prop:robust} (Direction-Corrected)}
Let $\Delta := g(x+\delta) - g(x)$ and let $S:=\Exp\Delta^2$. By
Section~4.1, $\Delta$ is a polynomial of degree at most $d^\star_{p99}$
in the Brownian increments of the integration grid, which are
independent standard Gaussians up to grid scaling.

Applying Theorem~\ref{thm:cw} with $p=\Delta$ and
$\varepsilon_{\mathrm{CW}} = \beta / S^{1/2}$ yields
\begin{equation}\label{eq:cw-exact}
  \Prob[|\Delta|\leq \beta] \;\leq\; C\,d^\star_{p99}\,
                                    \bigl(\beta/S^{1/2}\bigr)^{1/d^\star_{p99}}.
\end{equation}
Now let $L_{\mathrm{lo}}(\varepsilon)$ denote a Hoeffding lower
$(1-\alpha)$ confidence bound on $S^{1/2}/\varepsilon$ constructed from
$K\cdot N$ i.i.d.\ probe samples of the smoothed margin. Explicitly,
using the two-sample plug-in estimator $\hat S^{1/2}/\varepsilon$ of
the ratio and the fact that $\hat S/\varepsilon^2$ is bounded above by
$R^2$ under Assumption~\ref{ass:impl}, Hoeffding gives
\begin{equation*}
  \Prob\!\Bigl[\hat S^{1/2}/\varepsilon
                \geq S^{1/2}/\varepsilon - t\Bigr]
  \geq 1-\alpha, \quad
  t = R\sqrt{\log(1/\alpha)/(2KN)}.
\end{equation*}
Denote the confidence bound $L_{\mathrm{lo}}=\hat S^{1/2}/\varepsilon
- t$. With probability at least $1-\alpha$ we then have
$S^{1/2}\geq L_{\mathrm{lo}}\varepsilon$, and substituting into
\eqref{eq:cw-exact} preserves the direction of the inequality:
\begin{equation*}
  \Prob[|\Delta|\leq \beta] \;\leq\;
    C\,d^\star_{p99}\,\bigl(\beta/(L_{\mathrm{lo}}\varepsilon)\bigr)^{1/d^\star_{p99}},
\end{equation*}
which is exactly \eqref{eq:robust-bound}. \hfill$\square$

\section{Proof of Corollary~\ref{cor:argmax} (Decision-Flip Bound)}
A decision flip occurs at $x+\delta$ if and only if some competing
class $k\neq \arg\max g(x)$ satisfies $g_k(x+\delta)\geq g_{\arg\max}(x+\delta)$.
Let $m_k(x):=g_{\arg\max}(x)-g_k(x)$ denote the pairwise margin. The
flip event decomposes as
$\bigcup_{k\neq \arg\max}\{g_k(x+\delta)-g_{\arg\max}(x+\delta)\geq 0\}
 = \bigcup_{k}\{|\Delta_k|\geq m_k(x)\}$,
where $\Delta_k$ is the pairwise-margin gap. Applying
Proposition~\ref{prop:robust} to each $\Delta_k$ with $\beta=m_k(x)$
and a union bound over the $K-1$ classes yields
$\pi_{\mathrm{df}}(\varepsilon)\leq \sum_{k\neq \arg\max}
   C\,d^\star_{p99}\bigl(m_k(x)/(L_{\mathrm{lo}}\varepsilon)\bigr)^{1/d^\star_{p99}}
   \leq K\,C\,d^\star_{p99}
   \bigl(m(x)/(L_{\mathrm{lo}}\varepsilon)\bigr)^{1/d^\star_{p99}}$,
where $m(x)=\min_k m_k(x)$ is the confidence margin.
Inverting $\pi_{\mathrm{df}}(\varepsilon)\leq 1-c$ for $\varepsilon$
gives the radius \eqref{eq:radius}. \hfill$\square$

\section{BEL Bias / Variance Diagnostic}
Implementation in
\texttt{src/theory/bel\_estimator.py}. Protocol: pick
$\phi(X_T)=\|X_T\|^2$; compute a symmetric finite-difference reference
gradient of $\Exp\phi$ w.r.t.\ $\theta_g$ at step $h=10^{-3}$ (2048
paths per side); compute the BEL estimator on the same seeds with
$N\in\{32,128,512,2048\}$; report the mean and standard error of the
$\ell_2$ norm of the gradient. Results in Table~\ref{tab:bel-diag}.
The relative bias falls from 0.19 at $N=32$ to 0.023 at $N=2048$; the
standard error scales as $\mathcal O(N^{-1/2})$ as predicted by the
Central Limit Theorem.

\section{Stability Audit and Condition-Number Analysis}
Implementation in \texttt{src/theory/pseudoinverse.py}. Under
Assumption~\ref{ass:impl}, the condition number of
$g_\theta^\top g_\theta + \lambda_0 I_m$ has median $8.9$, 95-th
percentile $17.4$, worst-case $32.4$ across the ICS3D test partition.
Milstein comparison at $\Delta t / 2$ gives median terminal-state
discrepancy $1.4\times 10^{-4}$ (99-th percentile $5.3\times 10^{-4}$).
Spectral-norm audit: the largest singular value of each linear layer
in $f_\theta$, $g_\theta$ stays within $[0.86, 1.00]$ across the 40
training epochs.

\section{Dataset Provenance, Deduplication, and Leakage}
Full capture topology for ICS3D and IDS-PQC (tools, kernel versions,
capture windows, host and network configurations) is documented in
\texttt{docs/DATASETS.md}. Per-corpus unique-rate after BLAKE2b
fingerprinting: ICS3D $99.4\%$, IIS3D $97.8\%$, IDS-PQC $99.9\%$.
Cross-split fingerprint leakage under the temporal-holdout protocol:
$0.02\%$, $0.05\%$, $0.00\%$ respectively. IIS3D harmonisation
pipeline (schema alignment, missing-field imputation, label
canonicalisation) is a public Python module,
\texttt{src/data/harmonize/iis3d.py}, whose diff-based validation
against the shipped Kaggle CSVs is bit-exact.

%% ====================================================================
%%  ACKNOWLEDGEMENTS AND DATA AVAILABILITY
%% ====================================================================
\section*{Acknowledgments}
The authors thank the contributors to the open-source RobustIDPS.ai
platform and the maintainers of PyTorch and PyTorch Geometric, and the
anonymous reviewers for their constructive and detailed comments that
substantially strengthened the paper. This work was supported by a
grant for research centers in Artificial Intelligence provided by the
Ministry of Economic Development of the Russian Federation under
subsidy agreement (identifier 000000C313925P3Q0002). All authors of
this publication are co-first authors.

\section*{Data and Code Availability}
Stochastic ODE-Guard is implemented in a public reproducibility
repository\footnote{\url{https://github.com/rogerpanel/CV/tree/main/SODE-Guard}}
and as the \texttt{sode\_guard} detector within the open-source
RobustIDPS.ai platform.\footnote{\url{https://github.com/rogerpanel/robustidps.ai}}.
The three benchmark suites are released on Kaggle: ICS3D
(DOI \texttt{10.34740/kaggle/dsv/12483891}), IIS3D
(DOI \texttt{10.34740/kaggle/dsv/12479689}), and IDS-PQC
(DOI \texttt{10.34740/kaggle/dsv/15424420}), all under the CC BY-NC-SA
4.0 license.

%% ====================================================================
%% REFERENCES
%% ====================================================================
\bibliographystyle{IEEEtran}
\begin{thebibliography}{99}

\bibitem{wu2022rtids} Z.~Wu, H.~Zhang, P.~Wang, and Z.~Sun,
``RTIDS: A robust transformer-based approach for intrusion detection
system,'' \emph{IEEE Access}, vol.~10, pp.~64\,375--64\,387, 2022,
doi:~10.1109/ACCESS.2022.3182333.

\bibitem{halbouni2022cnnlstm} A.~Halbouni, T.~S.~Gunawan, M.~Habaebi,
M.~Halbouni, M.~Kartiwi, and R.~Ahmad, ``CNN-LSTM: Hybrid deep
neural network for network intrusion detection system,''
\emph{IEEE Access}, vol.~10, pp.~99\,837--99\,849, 2022,
doi:~10.1109/ACCESS.2022.3206425.

\bibitem{lo2022egraphsage} W.~W.~Lo, S.~Layeghy, M.~Sarhan, M.~Gallagher,
and M.~Portmann, ``E-GraphSAGE: A graph neural network based intrusion
detection system for IoT,'' in \emph{Proc.\ NOMS}, 2022, pp.~1--9,
doi:~10.1109/NOMS54207.2022.9789878; arXiv:2103.16329.

\bibitem{mosaiyebzadeh2025graphmamba} F.~Mosaiyebzadeh \textit{et~al.},
``IDS-GraphMamba: A Markov-enhanced graph Mamba framework for
real-time intrusion detection in IoMT edge networks,''
\emph{Computer Networks}, vol.~268, art.~108989, 2025,
doi:~10.1016/j.comnet.2025.108989.

\bibitem{ennaji2024advsurvey} S.~Ennaji \textit{et~al.}, ``Adversarial
challenges in network intrusion detection systems: Research insights
and future prospects,'' 2024, arXiv:2409.18736.

\bibitem{vitorino2024advbench} J.~Vitorino \textit{et~al.}, ``An
adversarial robustness benchmark for enterprise network intrusion
detection,'' 2024, arXiv:2402.16912.

\bibitem{venturi2024structural} A.~Venturi, D.~Stabili, and M.~Marchetti,
``Problem-space structural adversarial attacks for network intrusion
detection systems based on graph neural networks,'' 2024, arXiv:2403.11830.

\bibitem{tzen2019nsde} B.~Tzen and M.~Raginsky, ``Neural stochastic
differential equations: Deep latent Gaussian models in the diffusion
limit,'' 2019, arXiv:1905.09883.

\bibitem{li2020scalable} X.~Li, T.-K.~L.~Wong, R.~T.~Q.~Chen, and
D.~Duvenaud, ``Scalable gradients for stochastic differential
equations,'' in \emph{Proc.\ AISTATS}, PMLR vol.~108,
pp.~3870--3882, 2020; arXiv:2001.01328.

\bibitem{kidger2021nsdegan} P.~Kidger, J.~Foster, X.~Li, H.~Oberhauser,
and T.~Lyons, ``Neural SDEs as infinite-dimensional GANs,''
in \emph{Proc.\ ICML}, PMLR vol.~139, pp.~5453--5463, 2021;
arXiv:2102.03657.

\bibitem{oh2024stablesde} Y.~Oh, D.~Lim, and S.~Kim, ``Stable Neural
Stochastic Differential Equations in analyzing irregular time
series data,'' in \emph{Proc.\ ICLR}, 2024, arXiv:2402.14989.

\bibitem{kang2021sodef} Q.~Kang, Y.~Song, Q.~Ding, and W.~P.~Tay,
``Stable Neural ODE with Lyapunov-stable equilibrium points for
defending against adversarial attacks,'' in \emph{Proc.\ NeurIPS}, 2021,
arXiv:2110.12976.

\bibitem{guglielmi2025minimal} N.~Guglielmi \textit{et~al.},
``Minimal weight perturbation for robust Neural ODEs,'' 2025,
arXiv:2501.10740.

\bibitem{bergna2024lgnsde} R.~Bergna \textit{et~al.},
``Latent graph neural stochastic differential equations,'' 2024,
arXiv:2408.16115.

\bibitem{song2021score} Y.~Song \textit{et~al.}, ``Score-based
generative modeling through stochastic differential equations,''
in \emph{Proc.\ ICLR}, 2021, arXiv:2011.13456.

\bibitem{oksendal2003} B.~\O{}ksendal, \emph{Stochastic Differential
Equations: An Introduction with Applications}, 6th~ed. Springer, 2003,
doi:~10.1007/978-3-642-14394-6.

\bibitem{bismut1981martingales} J.-M.~Bismut, ``Martingales, the
Malliavin calculus and hypoellipticity under general H\"ormander's
conditions,'' \emph{Z.\ Wahrsch.\ Verw.\ Gebiete}, vol.~56, no.~4,
pp.~469--505, 1981, doi:~10.1007/BF00531428.

\bibitem{elworthy1994formulae} K.~D.~Elworthy and X.-M.~Li,
``Formulae for the derivatives of heat semigroups,''
\emph{J.\ Funct.\ Anal.}, vol.~125, no.~1, pp.~252--286, 1994,
doi:~10.1006/jfan.1994.1124.

\bibitem{hsu2002} E.~P.~Hsu, \emph{Stochastic Analysis on Manifolds},
Graduate Studies in Mathematics, vol.~38. AMS, 2002,
doi:~10.1090/gsm/038.

\bibitem{nualart2006malliavin} D.~Nualart, \emph{The Malliavin
Calculus and Related Topics}, 2nd~ed. Springer, 2006,
doi:~10.1007/3-540-28329-3.

\bibitem{carbery2001anticonc} A.~Carbery and J.~Wright,
``Distributional and $L^q$ norm inequalities for polynomials over
convex bodies in $\mathbb{R}^n$,'' \emph{Math.\ Res.\ Letters},
vol.~8, no.~3, pp.~233--248, 2001, doi:~10.4310/MRL.2001.v8.n3.a1.

\bibitem{meka2016anticonc} R.~Meka, O.~Nguyen, and V.~Vu,
``Anti-concentration for polynomials of independent random
variables,'' \emph{Theory of Computing}, vol.~12, art.~11, 2016,
doi:~10.4086/toc.2016.v012a011.

\bibitem{maurer2004pacbayes} A.~Maurer, ``A note on the PAC-Bayesian
theorem,'' 2004, arXiv:cs/0411099.

\bibitem{mcallester1999pacbayes} D.~A.~McAllester, ``PAC-Bayesian
model averaging,'' in \emph{Proc.\ COLT}, 1999, pp.~164--170.

\bibitem{dziugaite2017pacbayes} G.~K.~Dziugaite and D.~M.~Roy,
``Computing nonvacuous generalization bounds for deep (stochastic)
neural networks with many more parameters than training data,''
in \emph{Proc.\ UAI}, 2017, arXiv:1703.11008.

\bibitem{perezortiz2021pacbayes} M.~P\'erez-Ortiz, O.~Rivasplata,
J.~Shawe-Taylor, and Cs.~Szepesv\'ari, ``Tighter risk certificates for
neural networks,'' \emph{JMLR}, vol.~22, no.~227, pp.~1--40, 2021.

\bibitem{reeb2014klinv} D.~Reeb and Y.~Seldin, ``PAC-Bayes bounds for
supervised classification,'' 2014, arXiv:1410.0102.

\bibitem{cohen2019smoothing} J.~Cohen, E.~Rosenfeld, and J.~Z.~Kolter,
``Certified adversarial robustness via randomized smoothing,''
in \emph{Proc.\ ICML}, PMLR vol.~97, pp.~1310--1320, 2019,
arXiv:1902.02918.

\bibitem{salman2019smoothadv} H.~Salman \textit{et~al.}, ``Provably
robust deep learning via adversarially trained smoothed classifiers,''
in \emph{Proc.\ NeurIPS}, 2019, arXiv:1906.04584.

\bibitem{yang2020smoothing} G.~Yang \textit{et~al.}, ``Randomized
smoothing of all shapes and sizes,'' in \emph{Proc.\ ICML},
PMLR vol.~119, pp.~10\,693--10\,705, 2020, arXiv:2002.08118.

\bibitem{eiras2022ancer} F.~Eiras \textit{et~al.}, ``ANCER:
Anisotropic certification via sample-wise volume maximization,''
\emph{TMLR}, 2022, arXiv:2107.04570.

\bibitem{lecuyer2019pixeldp} M.~Lecuyer, V.~Atlidakis, R.~Geambasu,
D.~Hsu, and S.~Jana, ``Certified robustness to adversarial examples
with differential privacy,'' in \emph{Proc.\ IEEE S\&P}, 2019,
pp.~656--672, doi:~10.1109/SP.2019.00044.

\bibitem{tsuzuku2018lipschitz} Y.~Tsuzuku, I.~Sato, and M.~Sugiyama,
``Lipschitz-margin training: Scalable certification of perturbation
invariance for deep neural networks,'' in \emph{Proc.\ NeurIPS}, 2018,
arXiv:1802.04034.

\bibitem{anil2019sorting} C.~Anil, J.~Lucas, and R.~Grosse,
``Sorting out Lipschitz function approximation,'' in \emph{Proc.\ ICML},
PMLR vol.~97, pp.~291--301, 2019, arXiv:1811.05381.

\bibitem{wang2023lbdn} R.~Wang and I.~R.~Manchester, ``Direct
parameterization of Lipschitz-bounded deep networks,''
in \emph{Proc.\ ICML}, 2023, arXiv:2301.11526.

\bibitem{araujo2023sll} A.~Araujo \textit{et~al.}, ``A unified
algebraic perspective on Lipschitz neural networks,''
in \emph{Proc.\ ICLR}, 2023, arXiv:2303.03169.

\bibitem{mackay1992bnn} D.~J.~C.~MacKay, ``A practical Bayesian
framework for backpropagation networks,'' \emph{Neural Comput.},
vol.~4, no.~3, pp.~448--472, 1992.

\bibitem{liu2019advbnn} X.~Liu, Y.~Li, C.~Wu, and C.-J.~Hsieh,
``Adv-BNN: Improved adversarial defense through robust Bayesian
neural network,'' in \emph{Proc.\ ICLR}, 2019, arXiv:1810.01279.

\bibitem{he2019pni} Z.~He, A.~S.~Rakin, and D.~Fan,
``Parametric noise injection: Trainable randomness to improve deep
neural network robustness against adversarial attack,''
in \emph{Proc.\ CVPR}, 2019, pp.~588--597,
doi:~10.1109/CVPR.2019.00068.

\bibitem{snort3} The Snort Project, ``Snort~3 with SnortML neural
exploit-detection engine,'' Cisco Talos, Mar.~2024.
[Online]. Available: \url{https://github.com/snort3/snort3}.

\bibitem{suricata7} The Open Information Security Foundation,
``Suricata~7.0 release notes,'' 2024--2026.
[Online]. Available: \url{https://suricata.io/release-notes-7-0/}.

\bibitem{zeek7} The Zeek Project, ``Zeek~7.2 release,'' 2025.
[Online]. Available: \url{https://zeek.org/news/zeek-7-2}.

\bibitem{llamaguard3} Llama~Team, AI~@~Meta, ``Meta Llama Guard 3,''
\emph{Llama~3 technical report}, \S5.4, 2024, arXiv:2407.21783.

\bibitem{anaedevha2026mambashield} R.~N.~Anaedevha, A.~G.~Trofimov, and
Y.~V.~Borodachev, ``MambaShield: Selective state-space models for
poisoning-resilient sequential modeling,'' \emph{Expert Systems with
Applications}, Elsevier, 2026,
doi:~10.1016/j.eswa.2026.131175.

\bibitem{anaedevha2026uchgp} R.~N.~Anaedevha, A.~G.~Trofimov, and
Y.~V.~Borodachev, ``Uncertainty-calibrated hierarchical Gaussian
processes for intrusion detection with multi-scale temporal
modeling,'' \emph{Neurocomputing}, art.~133105, Elsevier, 2026,
doi:~10.1016/j.neucom.2026.133105.

\bibitem{anaedevha2026ode} R.~N.~Anaedevha, A.~G.~Trofimov, and
Y.~V.~Borodachev, ``Temporal adaptive neural ordinary differential
equations with deep spatio-temporal point processes for real-time
network intrusion detection,'' \emph{Complex \& Intelligent Systems},
Springer, 2026, doi:~10.1007/s40747-026-02291-7.

\bibitem{anaedevha2024ram} R.~N.~Anaedevha and A.~G.~Trofimov,
``Improved robust adversarial model against evasion attacks on
intrusion detection systems,'' \emph{Optical Memory and Neural
Networks (Information Optics)}, vol.~33, no.~S3, pp.~S414--S423,
Springer, 2024, doi:~10.3103/S1060992X24700681.

\bibitem{anaedevha2025stochtx} R.~N.~Anaedevha and A.~G.~Trofimov,
``Stochastic multimodal transformer with uncertainty quantification
for robust network intrusion detection,'' in \emph{Advances in Neural
Computation, Machine Learning, and Cognitive Research IX.
NEUROINFORMATICS 2025}, Studies in Computational Intelligence,
vol.~1241. Springer, 2025,
doi:~10.1007/978-3-032-07690-8\_35.

\bibitem{anaedevha2024elcon} R.~N.~Anaedevha, ``Application of machine
learning models for device identification in wireless network
traffic,'' in \emph{Proc.\ 2024 Conf.\ of Young Researchers in
Electrical and Electronic Engineering (ElCon)}, IEEE, 2024,
pp.~104--110, doi:~10.1109/ElCon61730.2024.10468413.

\bibitem{ds_ics3d} R.~N.~Anaedevha, ``Integrated Cloud Security
3-Datasets (ICS3D),'' Kaggle, 2025,
doi:~10.34740/kaggle/dsv/12483891.

\bibitem{ds_iis3d} R.~N.~Anaedevha, ``Integrated IDPS Security
3-Datasets (IIS3D),'' Kaggle, 2025,
doi:~10.34740/kaggle/dsv/12479689.

\bibitem{ds_pqcids} R.~N.~Anaedevha, ``IDS encrypted and
post-quantum-cryptography datasets,'' Kaggle, 2026,
doi:~10.34740/kaggle/dsv/15424420.

\bibitem{robustidps_repo} R.~N.~Anaedevha,
``RobustIDPS.ai v1.1.0 open-source AI-native IDPS,''
GitHub repository, MIT License, 2026.
[Online]. Available: \url{https://github.com/rogerpanel/robustidps.ai}.

\end{thebibliography}

\vspace{-1.0cm}

\begin{IEEEbiography}[{\includegraphics[width=1in,height=1.25in,clip,keepaspectratio]{author1.jpeg}}]{Roger Nick Anaedevha}
received the B.Sc. degree in Computer Science from Ambrose Alli University, Nigeria, M.Sc. degrees in Computer Science from the University of Abuja, Nigeria, and in Information Security from the National Research Nuclear University MEPhI, Moscow. He is currently pursuing the Ph.D. degree in Applied Mathematics and Artificial Intelligence at MEPhI under Prof. Trofimov Aleksandr Gennadievich. His research interests include adversarial machine learning, federated learning, privacy-preserving AI, uncertainty quantification, and Quantum Machine Learning. He has authored several first-author publications in IEEE, Springer, and Elsevier venues. His adversarially-robust IDPS work has been deployed in production processing over 2\,TB of daily traffic. He is a member of IEEE, ACM, and Russian Neural Network Association.
\end{IEEEbiography}

\vspace{-1.0cm}

\begin{IEEEbiography}[{\includegraphics[width=1in,height=1.25in,clip,keepaspectratio]{author2.jpeg}}]{Trofimov Aleksandr Gennadievich}
has a PhD, the Assistant Professor at the Institute of Cyber Intelligent Systems, National Research Nuclear University, MEPhI. Head of the Laboratory of Neural Networks Technologies (MEPhI), program/organizing committee member of International Conference ``Neuroinformatics'', the research advisor at the Center for Top-level educational programs in the field of AI at the Institute of Cyber Intelligent Systems of MEPhI. His research focuses on intelligent systems, machine learning theory, and security applications.
\end{IEEEbiography}

\vspace{-1.0cm}

\begin{IEEEbiography}[{\includegraphics[width=1in,height=1.25in,clip,keepaspectratio]{author3.png}}]{Yuri Vladimirovich Borodachev}
is affiliated with the Artificial Intelligence Research Center at the National Research Nuclear University MEPhI, Moscow, Russia. He is responsible for managing, coordinating and organizing the research work. His research interests include engineering systems based on artificial intelligence, autonomous vehicles, and AI applications in transport.
\end{IEEEbiography}

\end{document}
